% ============================================================
%  Videogame Semantic Search — Presentazione Beamer
%  Focus: l'ontologia OWL 2 DL (14 slide di contenuto)
% ============================================================
\documentclass[aspectratio=169,11pt]{beamer}
\usetheme[progressbar=frametitle, numbering=fraction]{metropolis}

% --- Encoding e lingua ---
\usepackage[italian]{babel}
\usepackage[utf8]{inputenc}
\usepackage[T1]{fontenc}

% --- Tabelle ---
\usepackage{booktabs}
\usepackage{tabularx}
\usepackage{array}

% --- Codice ---
\usepackage{listings}
\usepackage{xcolor}

% --- Grafica ---
\usepackage{tikz}
\usetikzlibrary{positioning, arrows.meta, shapes.geometric}

% --- Matematica ---
\usepackage{amsmath, amssymb}

% --- Nascondi simboli navigazione ---
\setbeamertemplate{navigation symbols}{}

% --- Margini ---
\setbeamersize{text margin left=10mm, text margin right=10mm}

% --- Colori accent ---
\definecolor{accent}{HTML}{3B82F6}
\setbeamercolor{progress bar}{fg=accent}
\setbeamercolor{frametitle}{bg=accent}
\setbeamercolor{alerted text}{fg=accent}

% --- Stile listings (sobrio) ---
\definecolor{codebg}{rgb}{0.94,0.94,0.94}
\definecolor{codecomment}{rgb}{0.4,0.4,0.4}
\definecolor{codekeyword}{rgb}{0.0,0.0,0.6}
\definecolor{codestring}{rgb}{0.5,0.1,0.1}

\lstdefinestyle{sparql}{
  backgroundcolor=\color{codebg},
  basicstyle=\ttfamily\scriptsize,
  keywordstyle=\color{codekeyword}\bfseries,
  commentstyle=\color{codecomment}\itshape,
  stringstyle=\color{codestring},
  breaklines=true,
  frame=single,
  rulecolor=\color{codecomment},
  numbers=none,
  showstringspaces=false,
  tabsize=2,
  aboveskip=2pt,
  belowskip=2pt,
  morestring=[b]",
  morecomment=[l]\#,
  morekeywords={PREFIX,SELECT,DISTINCT,WHERE,FILTER,ORDER,BY,DESC,LIMIT,a},
  columns=fullflexible,
  keepspaces=true
}

\lstdefinestyle{turtle}{
  backgroundcolor=\color{codebg},
  basicstyle=\ttfamily\scriptsize,
  keywordstyle=\color{codekeyword}\bfseries,
  commentstyle=\color{codecomment}\itshape,
  stringstyle=\color{codestring},
  breaklines=true,
  frame=single,
  rulecolor=\color{codecomment},
  numbers=none,
  showstringspaces=false,
  tabsize=2,
  aboveskip=2pt,
  belowskip=2pt,
  morestring=[b]",
  morecomment=[l]\#,
  morekeywords={a},
  columns=fullflexible,
  keepspaces=true
}

% --- Metadati -------------------------------------------------
\title{Videogame Semantic Search}
\subtitle{Un'ontologia OWL 2 DL per il dominio dei videogiochi \\ Modellazione, reasoning e interrogazione}
\author{Videogame Semantic Search Project}
\date{Agosto 2026}

% ============================================================
\begin{document}

% ============================================================
%  SLIDE 0 — Titolo
% ============================================================
\begin{frame}[plain]
  \titlepage
\end{frame}

% ============================================================
%  SLIDE 1 — Obiettivo
% ============================================================
\begin{frame}{Obiettivo del Progetto}
  \begin{block}{Cosa abbiamo costruito}
    Un'\textbf{ontologia OWL 2 DL} sul dominio dei videogiochi (2010--2026),
    popolata con dati estratti da Wikidata e interrogabile tramite SPARQL.
  \end{block}

  \vspace{4pt}

  \begin{itemize}
    \item \textbf{Perch\'e}: rappresentare in modo formale e ragionabile la conoscenza su
      giochi, studi, piattaforme, generi, franchise e premi
    \item \textbf{Cosa contiene}:
    \begin{itemize}
      \item 78 classi (15 definite) e 57 propriet\`a con caratteristiche OWL 2
      \item 12 blocchi di disjointness, 3 property chain, 3 chiavi (\texttt{hasKey})
      \item Allineamento a Wikidata e Dublin Core
    \end{itemize}
    \item \textbf{Come si usa}: reasoning a due livelli (OWL 2 RL + SPARQL CONSTRUCT) che
      classifica automaticamente le istanze; un'applicazione web (NL $\to$ SPARQL, grafo
      interattivo) la rende esplorabile
  \end{itemize}
\end{frame}

% ============================================================
%  SLIDE 2 — Panoramica
% ============================================================
\begin{frame}{Ontologia --- Panoramica}
  \begin{itemize}\small
    \item \textbf{Dominio}: videogiochi 2010--2026 con sviluppatori, publisher, generi,
      piattaforme, personaggi, franchise, premi, motori, modalit\`a, recensioni ed eventi
    \item \textbf{Namespace}: \texttt{http://www.videogame-ontology.org/ontology\#} (prefisso \texttt{vg:})
    \item \textbf{Schema e istanze separati}: \texttt{videogames.owl} (schema, 934 triple) e
      \texttt{videogames\_pruned\_2020.owl} (istanze demo, $\sim$745k triple)
  \end{itemize}

  \vspace{2pt}

  \begin{table}
  \centering\small
  \begin{tabular}{@{}lrlr@{}}
    \toprule
    \textbf{Componente} & \textbf{Qt\`a} & \textbf{Componente} & \textbf{Qt\`a} \\
    \midrule
    Classi totali            & 78  & Classi definite (\texttt{equivalentClass}) & 15 \\
    Classi primitive         & 61  & Classi con vincolo di cardinalit\`a         & 2  \\
    Object property          & 31  & Blocchi \texttt{AllDisjointClasses}         & 12 \\
    Datatype property        & 26  & \texttt{propertyChainAxiom}                 & 3  \\
    Coppie inverse           & 12  & \texttt{hasKey}                             & 3  \\
    Propriet\`a transitive   & 3   & Allineamenti (Wikidata + Dublin Core)       & 4 + 3 \\
    \bottomrule
  \end{tabular}
  \end{table}
\end{frame}

% ============================================================
%  SLIDE 3 — Gerarchia classi
% ============================================================
\begin{frame}{Ontologia --- Gerarchia delle Classi}
  \begin{columns}[T]
    \column{0.44\textwidth}
    \textbf{12 classi top-level} {\scriptsize (disgiunte)}
    \vspace{2pt}

    \small
    \begin{columns}[T]
      \column{0.5\linewidth}
      \begin{itemize}\setlength{\itemsep}{-1pt}
        \item VideoGame
        \item Developer
        \item Publisher
        \item Genre
        \item Platform
        \item Character
      \end{itemize}
      \column{0.5\linewidth}
      \begin{itemize}\setlength{\itemsep}{-1pt}
        \item Franchise
        \item Award
        \item GameEngine
        \item GameMode
        \item Review
        \item GameEvent
      \end{itemize}
    \end{columns}

    \column{0.54\textwidth}
    \textbf{Platform --- partizionamento a due livelli}
    \vspace{2pt}

    \small
    \begin{itemize}
      \item \textbf{ConsolePlatform}: PlayStation, Xbox, Nintendo
      \item \textbf{PCPlatform}: Windows, MacOS, Linux
      \item \textbf{MobilePlatform}: iOS, Android
    \end{itemize}
    \vspace{2pt}
    Altre sottoclassi: 10 generi (aperti), stato di rilascio, budget (Indie/AAA),
    tipologie di premio, personaggio e motore
  \end{columns}

  \vspace{6pt}

  \begin{block}{Pattern modellistici}
    \small
    \textbf{Review} --- relazione N-aria: gioco (\texttt{reviewedGame}) + revisore + punteggio + data \\
    \textbf{GameEvent} --- ciclo di vita: Announcement, Release, Update, Delisting \\
    \textbf{GameMode} --- value partition: 4 modalit\`a disgiunte
  \end{block}
\end{frame}

% ============================================================
%  SLIDE 4 — Proprietà: caratteristiche
% ============================================================
\begin{frame}{Ontologia --- Caratteristiche delle Propriet\`a}
  \footnotesize
  \begin{tabularx}{\textwidth}{@{}l r >{\raggedright\arraybackslash}X@{}}
    \toprule
    \textbf{Caratteristica} & \textbf{\#} & \textbf{Esempi} \\
    \midrule
    \texttt{owl:inverseOf}          & 12 coppie & \texttt{developedBy}/\texttt{developerOf}, \texttt{belongsTo}/\texttt{includes}, \dots \\
    \texttt{TransitiveProperty}     & 3  & \texttt{sequelOf}, \texttt{partOfFranchise}, \texttt{subsidiaryOf} \\
    \texttt{SymmetricProperty}      & 3  & le super-propriet\`a delle 3 catene (\texttt{sharesDeveloperWith}, \dots) \\
    \texttt{Asymmetric} / \texttt{Irreflexive} & 13 / 18 & \texttt{developedBy}, \texttt{hasGenre}, \texttt{belongsTo}, \dots \\
    \texttt{FunctionalProperty}     & 14 & \texttt{reviewedGame} + 13 datatype (\texttt{releaseDate}, \texttt{metacriticScore}, \dots) \\
    \texttt{subPropertyOf}          & 9  & \texttt{gameName}, \texttt{developerName}, \dots{} $\sqsubseteq$ \texttt{name} \\
    \bottomrule
  \end{tabularx}

  \vspace{4pt}

  \begin{alertblock}{Vincolo globale di OWL 2 DL}
    \small
    Le transitive e le loro inverse non sono \emph{semplici}: non possono essere asimmetriche,
    irriflessive, funzionali, n\'e comparire in cardinalit\`a o \texttt{hasKey}.
  \end{alertblock}
\end{frame}

% ============================================================
%  SLIDE 5 — Property chain e hasKey
% ============================================================
\begin{frame}{Ontologia --- Property Chain e Chiavi}
  \begin{block}{3 Property Chain Axiom (inclusione di ruoli)}
    \small
    \vspace{-4pt}
    \begin{align*}
      \texttt{belongsTo} \circ \texttt{includes} &\sqsubseteq \texttt{sharedFranchiseWith} \\
      \texttt{developedBy} \circ \texttt{developerOf} &\sqsubseteq \texttt{sharesDeveloperWith} \\
      \texttt{publishedBy} \circ \texttt{publisherOf} &\sqsubseteq \texttt{sharesPublisherWith}
    \end{align*}
  \end{block}

  \begin{itemize}\small
    \item Composizione ben tipata: VideoGame $\xrightarrow{\texttt{belongsTo}}$ Franchise
      $\xrightarrow{\texttt{includes}}$ VideoGame
    \item Le super-propriet\`a sono \texttt{Symmetric}; la catena include anche la coppia
      $(g,g)$, per questo nelle query si filtra \texttt{?g != ?altro}
  \end{itemize}

  \begin{block}{3 chiavi \texttt{owl:hasKey}}
    \footnotesize
    \texttt{VideoGame(hasWikidataId)} $\cdot$
    \texttt{VideoGame(hasSteamAppId)} $\cdot$
    \texttt{Review(reviewedGame, reviewerName)} \\[2pt]
    Identit\`a espressa con chiavi (valide sui soli individui nominati) anzich\'e con
    \texttt{InverseFunctionalProperty} su datatype, non ammessa in OWL 2 DL.
  \end{block}
\end{frame}

% ============================================================
%  SLIDE 6 — Classi definite (∃)
% ============================================================
\begin{frame}{Ontologia --- Classi Definite (Condizioni Necessarie e Sufficienti)}
  \small
  \begin{tabular}{@{}ll@{}}
    \toprule
    \textbf{Classe} & \textbf{Definizione (\texttt{owl:equivalentClass})} \\
    \midrule
    AwardWinningGame     & $\texttt{VideoGame} \sqcap \exists\,\texttt{wonAward.Award}$ \\
    FranchiseGame        & $\texttt{VideoGame} \sqcap \exists\,\texttt{belongsTo.Franchise}$ \\
    SequelGame           & $\texttt{VideoGame} \sqcap \exists\,\texttt{sequelOf.VideoGame}$ \\
    Multiplayer / SinglePlayer / CoopGame & $\texttt{VideoGame} \sqcap \exists\,\texttt{hasGameMode}.M$ \\
    PC / Console / MobileGame & $\texttt{VideoGame} \sqcap \exists\,\texttt{availableOn}.P$ \\
    PlayStation / Xbox / NintendoGame & $\texttt{VideoGame} \sqcap \exists\,\texttt{availableOn}.P$ \\
    \bottomrule
  \end{tabular}

  \vspace{2pt}
  {\footnotesize $M$ = la modalit\`a corrispondente (es.\ \texttt{MultiplayerMode});
  $P$ = la classe di piattaforma corrispondente (es.\ \texttt{PCPlatform}, \texttt{PlayStationPlatform})\par}

  \vspace{6pt}

  \begin{itemize}
    \item 12 classi definite con \texttt{intersectionOf} + \texttt{someValuesFrom} su object property
    \item Altre 3 (HighlyRated, Recent, Standalone) usano costrutti pi\`u espressivi $\to$ 15 in totale
    \item La gerarchia delle piattaforme propaga la classificazione
      (PS5 $\Rightarrow$ PlayStationGame e ConsoleGame)
  \end{itemize}
\end{frame}

% ============================================================
%  SLIDE 7 — Classi definite: altri costrutti
% ============================================================
\begin{frame}{Ontologia --- Negazione, Datatype e Cardinalit\`a}
  \footnotesize
  \begin{tabular}{@{}l l l@{}}
    \toprule
    \textbf{Classe} & \textbf{Assioma} & \textbf{Costrutto} \\
    \midrule
    StandaloneGame  & $\equiv \texttt{VideoGame} \sqcap \lnot\exists\,\texttt{sequelOf.VideoGame}$ & \texttt{complementOf} \\
    HighlyRatedGame & $\equiv \texttt{VideoGame} \sqcap \exists\,\texttt{metacriticScore}.\texttt{integer}[\ge 85]$ & \texttt{minInclusive} \\
    RecentGame      & $\equiv \texttt{VideoGame} \sqcap \exists\,\texttt{releaseDate}.\texttt{date}[\ge 2020\text{-}01\text{-}01]$ & \texttt{minInclusive} \\
    \midrule
    GameWithSingleDeveloper & $\sqsubseteq \texttt{VideoGame} \sqcap {\le}1\,\texttt{developedBy}$ & \texttt{maxCardinality} \\
    ExclusiveRelease        & $\sqsubseteq \texttt{VideoGame} \sqcap {\le}1\,\texttt{availableOn}$ & \texttt{maxCardinality} \\
    \bottomrule
  \end{tabular}

  {\scriptsize Le cardinalit\`a sono \emph{non qualificate} (nessun \texttt{owl:onClass}).\par}

  \vspace{6pt}

  \begin{itemize}
    \item Le ultime due hanno \textbf{sole condizioni necessarie}: per l'Open World
      Assumption un reasoner dedurrebbe ${\le}1$ solo con assiomi di chiusura, quindi fungono da
      vincoli di integrit\`a e non da classificatori
    \item \texttt{complementOf} e le datatype restriction \textbf{escono dal profilo OWL 2 RL}:
      motivano il secondo livello di reasoning
  \end{itemize}
\end{frame}

% ============================================================
%  SLIDE 8 — L'OWL in pratica
% ============================================================
\begin{frame}[fragile]{Ontologia --- Gli Assiomi in Turtle}
  {\small Estratti da \texttt{videogames.owl} (serializzati in Turtle)}
  \vspace{2pt}

  \begin{columns}[T]
    \column{0.50\textwidth}
\begin{lstlisting}[style=turtle]
# Classe definita con datatype restriction
vg:HighlyRatedGame a owl:Class ;
  owl:equivalentClass [
    owl:intersectionOf ( vg:VideoGame
      [ a owl:Restriction ;
        owl:onProperty vg:metacriticScore ;
        owl:someValuesFrom [
          a rdfs:Datatype ;
          owl:onDatatype xsd:integer ;
          owl:withRestrictions
            ( [ xsd:minInclusive 85 ] ) ] ] ) ] .
\end{lstlisting}

    \column{0.50\textwidth}
\begin{lstlisting}[style=turtle]
# Property chain + simmetria
vg:sharesDeveloperWith
  a owl:ObjectProperty ,
    owl:SymmetricProperty ;
  owl:propertyChainAxiom
    ( vg:developedBy vg:developerOf ) .

# Condizione solo necessaria
vg:GameWithSingleDeveloper
  rdfs:subClassOf vg:VideoGame ,
    [ a owl:Restriction ;
      owl:onProperty vg:developedBy ;
      owl:maxCardinality
        "1"^^xsd:nonNegativeInteger ] .
\end{lstlisting}
  \end{columns}
\end{frame}

% ============================================================
%  SLIDE 9 — Disjointness
% ============================================================
\begin{frame}{Ontologia --- Disjointness}
  \textbf{12 blocchi \texttt{owl:AllDisjointClasses}}: nessuna istanza pu\`o appartenere
  a due classi dello stesso blocco
  \vspace{4pt}

  \scriptsize
  \begin{columns}[T]
    \column{0.44\textwidth}
    \begin{tabular}{@{}r l@{}}
      \toprule
      \textbf{\#} & \textbf{Classi disgiunte} \\
      \midrule
      12 & le classi top-level \\
      3  & Console / PC / MobilePlatform \\
      3  & PlayStation / Xbox / NintendoPlatform \\
      3  & Windows / MacOS / LinuxPlatform \\
      2  & IOS / AndroidPlatform \\
      3  & Industry / Editorial / UserAward \\
      \bottomrule
    \end{tabular}

    \column{0.54\textwidth}
    \begin{tabular}{@{}r l@{}}
      \toprule
      \textbf{\#} & \textbf{Classi disgiunte} \\
      \midrule
      2 & Player / NonPlayerCharacter \\
      2 & Proprietary / OpenSourceEngine \\
      4 & Released / EarlyAccess / Cancelled / DelistedGame \\
      2 & IndieGame / AAAGame \\
      4 & SinglePlayer / Multiplayer / Coop / MMOMode \\
      4 & Announcement / Release / Update / DelistingEvent \\
      \bottomrule
    \end{tabular}
  \end{columns}

  \vspace{6pt}

  \small
  \begin{itemize}
    \item \textbf{Scelta deliberata}: SinglePlayerGame e MultiplayerGame \emph{non} sono disgiunte
      (un gioco pu\`o avere entrambe le modalit\`a); i generi restano una tassonomia aperta
    \item Con un reasoner DL le disjointness rilevano dati incoerenti provenienti da Wikidata
  \end{itemize}
\end{frame}

% ============================================================
%  SLIDE 10 — Allineamento e profilo
% ============================================================
\begin{frame}{Ontologia --- Allineamento Esterno e Profilo OWL 2 DL}
  \begin{columns}[T]
    \column{0.50\textwidth}
    \begin{block}{4 \texttt{equivalentClass} con Wikidata}
      \small
      \texttt{vg:VideoGame} $\equiv$ \texttt{wd:Q7889} \\
      \texttt{vg:Developer} $\equiv$ \texttt{wd:Q210167} \\
      \texttt{vg:GameEngine} $\equiv$ \texttt{wd:Q193564} \\
      \texttt{vg:Genre} $\equiv$ \texttt{wd:Q65956376}
    \end{block}

    \begin{block}{3 \texttt{equivalentProperty} con Dublin Core}
      \small
      \texttt{vg:releaseDate} $\equiv$ \texttt{dcterms:issued} \\
      \texttt{vg:gameName} $\equiv$ \texttt{dcterms:title} \\
      \texttt{vg:gameDescription} $\equiv$ \texttt{dcterms:description}
    \end{block}

    \column{0.48\textwidth}
    \begin{block}{Vincoli globali OWL 2 DL sulle propriet\`a}
      \small
      \begin{itemize}
        \item Nessuna propriet\`a non semplice in asimmetria, irriflessivit\`a,
          funzionalit\`a, cardinalit\`a o \texttt{hasKey}
        \item Nessuna \texttt{InverseFunctionalProperty} su datatype property
        \item Verifica: \textbf{0 violazioni}
      \end{itemize}
    \end{block}
    {\footnotesize Nota: \texttt{xsd:date} non \`e nella datatype map di OWL 2 (solo
    \texttt{xsd:dateTime}); per usare HermiT va migrato a \texttt{xsd:dateTime}.\par}
  \end{columns}
\end{frame}

% ============================================================
%  SLIDE 11 — Reasoning a due livelli
% ============================================================
\begin{frame}{Reasoning --- Architettura a Due Livelli}
  \begin{center}
  \begin{tikzpicture}[
    node distance=5mm,
    box/.style={rectangle, draw, rounded corners, minimum width=24mm, minimum height=11mm,
                fill=accent!10, font=\footnotesize, align=center},
    data/.style={rectangle, draw=gray, dashed, rounded corners, minimum width=18mm,
                 minimum height=9mm, font=\scriptsize, align=center},
    arr/.style={-{Stealth}, thick, accent}
  ]
    \node[data] (in) {Grafo RDF\\$\sim$997 triple};
    \node[box, right=of in] (l1) {\textbf{Livello 1}\\owlrl (OWL 2 RL)};
    \node[data, right=of l1] (mid) {$\sim$2.680\\triple};
    \node[box, right=of mid] (l2) {\textbf{Livello 2}\\21 regole CONSTRUCT};
    \node[data, right=of l2] (out) {Grafo inferito\\2.715 triple};
    \draw[arr] (in) -- (l1);
    \draw[arr] (l1) -- (mid);
    \draw[arr] (mid) -- (l2);
    \draw[arr] (l2) -- (out);
  \end{tikzpicture}
  \end{center}
  {\centering\scriptsize (valori sull'insieme di istanze di test: 3 giochi, 2 developer, piattaforme)\par}

  \vspace{4pt}

  \begin{columns}[T]
    \column{0.49\textwidth}
    \textbf{owlrl materializza}
    \begin{itemize}\small
      \item \texttt{subClassOf}, \texttt{intersectionOf} + \texttt{someValuesFrom}
      \item \texttt{inverseOf}, \texttt{Symmetric}, \texttt{Transitive}
      \item \texttt{propertyChainAxiom}
    \end{itemize}

    \column{0.49\textwidth}
    \textbf{Livello 2: regole CONSTRUCT}
    \begin{itemize}\small
      \item datatype restriction: \texttt{FILTER(?s >= 85)}
      \item \texttt{complementOf}: \texttt{FILTER NOT EXISTS} (NAF)
      \item classificazioni e inverse (ridondanti con RL)
    \end{itemize}
  \end{columns}

  \vspace{4pt}
  {\small ${\le}1$ in superclasse \`e ammessa in RL (\texttt{cls-maxc2}) ma non classifica: serve solo come vincolo.}
\end{frame}

% ============================================================
%  SLIDE 12 — Esempio d'inferenza
% ============================================================
\begin{frame}[fragile]{Reasoning --- Un Esempio di Inferenza}
  \begin{columns}[T]
    \column{0.40\textwidth}
    \textbf{Asserito}
\begin{lstlisting}[style=turtle]
# GameA = Elden Ring
vg:GameA a vg:VideoGame ;
  vg:releaseDate
      "2022-02-25"^^xsd:date ;
  vg:metacriticScore 96 ;
  vg:developedBy vg:Dev1 ;
  vg:hasGameMode vg:SinglePlayer ,
      vg:OnlineMultiplayer ;
  vg:availableOn vg:PS5 ,
      vg:PC_Windows ;
  vg:wonAward vg:TGAAward2022 .
vg:GameC vg:developedBy vg:Dev1 .
vg:PS5 a vg:PlayStationPlatform .
vg:PC_Windows a vg:WindowsPlatform .
\end{lstlisting}

    \column{0.58\textwidth}
    \textbf{Inferito}
    \vspace{2pt}

    \scriptsize
    \setlength{\tabcolsep}{2.5pt}
    \begin{tabular}{@{}l l@{}}
      \toprule
      \textbf{Tipi inferiti di GameA} & \textbf{Motivo} \\
      \midrule
      PlayStation-, ConsoleGame   & PS5 : PlayStationPlatform \\
      PCGame                      & Windows $\sqsubseteq$ PCPlatform \\
      Single-, MultiplayerGame    & $\exists$\,hasGameMode.$M$ \\
      AwardWinningGame            & $\exists$\,wonAward.Award \\
      \textit{HighlyRatedGame}    & \textit{$96 \ge 85$} \\
      \textit{RecentGame}         & \textit{$2022 \ge 2020$} \\
      \textit{StandaloneGame}     & \textit{nessun sequelOf} \\
      \midrule
      \textbf{Relazioni inferite} & \\
      Dev1 developerOf GameA      & inverseOf \\
      GameA sharesDeveloperWith GameC & catena + simm. \\
      \bottomrule
    \end{tabular}

    \vspace{3pt}
    {\scriptsize In tondo: livello 1 (owlrl). In corsivo: livello 2 (CONSTRUCT).\par}
  \end{columns}
\end{frame}

% ============================================================
%  SLIDE 13 — Query SPARQL
% ============================================================
\begin{frame}[fragile]{Interrogare l'Ontologia --- SPARQL}
  \begin{columns}[T]
    \column{0.50\textwidth}
    \textbf{Classi inferite al posto dei filtri}
\begin{lstlisting}[style=sparql]
PREFIX vg: <http://www.videogame-ontology.org/ontology#>
SELECT ?name ?score WHERE {
  ?g a vg:HighlyRatedGame ,
       vg:MultiplayerGame ,
       vg:PlayStationGame ;
     vg:gameName ?name ;
     vg:metacriticScore ?score .
}
ORDER BY DESC(?score) LIMIT 10
\end{lstlisting}
    {\scriptsize Nessun filtro su punteggio, modalit\`a o gerarchia delle piattaforme:
    li ha gi\`a risolti il reasoner.}

    \column{0.50\textwidth}
    \textbf{Relazioni derivate dalle catene}
\begin{lstlisting}[style=sparql]
PREFIX vg: <http://www.videogame-ontology.org/ontology#>
SELECT DISTINCT ?other WHERE {
  ?er vg:gameName "Elden Ring" ;
      vg:sharesDeveloperWith ?g .
  ?g  vg:gameName ?other .
  FILTER(?g != ?er)
}
\end{lstlisting}
    {\scriptsize \texttt{sharesDeveloperWith} non \`e mai asserita: deriva da
    \texttt{developedBy} $\circ$ \texttt{developerOf}. Il \texttt{FILTER} esclude la coppia
    riflessiva prodotta dalla catena.}
  \end{columns}
\end{frame}

% ============================================================
%  SLIDE 14 — Popolamento e architettura
% ============================================================
\begin{frame}{Popolamento e Contesto Applicativo}
  \begin{columns}[T]
    \column{0.52\textwidth}
    \textbf{Popolamento da Wikidata}
    \begin{enumerate}\small
      \item 12 tipologie di query SPARQL, ciascuna su 204 finestre anno $\times$ mese
      \item Deduplicazione: 4.420 entit\`a rimosse
      \item Chiusura OWL 2 RL: $\sim$1M $\to$ $\sim$1.4M triple
      \item Pruning: $\sim$1.17M triple (2010--2026)
      \item Demo 2020--2026: $\sim$745k triple, $\sim$68.7k giochi
    \end{enumerate}

    \column{0.46\textwidth}
    \textbf{Architettura applicativa}
    \vspace{2pt}

    \begin{tikzpicture}[
      node distance=3mm,
      box/.style={rectangle, draw, rounded corners, text width=54mm, minimum height=7mm,
                  fill=accent!10, font=\scriptsize, align=center},
      arr/.style={-{Stealth[scale=0.7]}, thick, accent}
    ]
      \node[box] (fe) {\textbf{Frontend} React + grafo interattivo};
      \node[box, below=of fe] (be) {\textbf{Backend} FastAPI + agente NL $\to$ SPARQL};
      \node[box, below=of be] (st) {\textbf{Store} pyoxigraph (Rust) sull'ontologia};
      \draw[arr] (fe) -- (be);
      \draw[arr] (be) -- (st);
    \end{tikzpicture}

    \vspace{4pt}
    {\scriptsize pyoxigraph vs rdflib: query 5--62$\times$ pi\`u veloci, caricamento
    $\sim$2\,s invece di $\sim$41\,s, RAM ridotta di circa 4$\times$.}
  \end{columns}
\end{frame}

% ============================================================
%  SLIDE FINALE — Thanks
% ============================================================
\begin{frame}[standout]
  \Huge Grazie per l'attenzione \\[12pt]
  \normalsize Videogame Semantic Search Project \\[6pt]
  \footnotesize Ontologia OWL 2 DL · Reasoning · Interrogazione SPARQL
\end{frame}

\end{document}
